\textbf{1. Kiến thức cần nhớ}
\begin{itemize}
\item \textbf{Đa thức} là một tổng của những đơn thức. Mỗi đơn thức trong tổng gọi là một \textit{hạng tử} của đa thức đó.
\begin{itemize}
\item Mỗi đơn thức cũng được coi là một đa thức (đa thức chỉ có một hạng tử).
\item Số /bin/zsh$ được gọi là \textit{đa thức không}.
\end{itemize}
\item \textbf{Đa thức thu gọn} là đa thức không còn hai hạng tử nào đồng dạng.
\item \textbf{Bậc của đa thức}:
\begin{itemize}
\item Bậc của đa thức là \textbf{bậc của hạng tử có bậc cao nhất} trong dạng thu gọn của đa thức đó.
\item Đa thức không không có bậc.
\item Một số thực khác /bin/zsh$ là đa thức bậc /bin/zsh$.
\end{itemize}
\item \textbf{Giá trị của đa thức}: Để tính giá trị của một đa thức tại các giá trị cho trước của biến, ta thay các giá trị đó vào đa thức rồi thực hiện phép tính.
\end{itemize}

\textbf{2. Phương pháp giải các dạng}

%% DANG 8D1?2-1
\textbf{Dạng 1. Nhận biết đa thức, thu gọn đa thức}
\begin{enumerate}
\item \textbf{Nhận biết}: Kiểm tra xem biểu thức có phải là tổng của các đơn thức hay không.
\item \textbf{Thu gọn đa thức}:
\begin{itemize}
\item Nhóm các hạng tử đồng dạng lại với nhau (đặt dấu cộng trước các ngoặc nhóm để tránh nhầm dấu).
\item Thực hiện phép cộng, trừ các đơn thức đồng dạng trong từng nhóm.
\end{itemize}
\end{enumerate}
Lưu ý: Khi nhóm các hạng tử mang dấu trừ, nên giữ nguyên dấu của hạng tử bên trong ngoặc, ví dụ: 569X3x^2y + 5x^2y = (-3 + 5)x^2y = 2x^2y$.

%% DANG 8D1?2-2
\textbf{Dạng 2. Bậc của đa thức}
\begin{enumerate}
\item \textbf{Bước 1}: \textbf{Bắt buộc thu gọn đa thức} trước khi xác định bậc (các hạng tử bậc cao có thể triệt tiêu lẫn nhau).
\item \textbf{Bước 2}: Tìm bậc của từng hạng tử trong đa thức đã thu gọn.
\item \textbf{Bước 3}: Chọn số lớn nhất trong các bậc đó làm bậc của đa thức.
\end{enumerate}
Sai lầm thường gặp: Quên không thu gọn đa thức dẫn đến xác định sai bậc khi hai hạng tử có bậc cao nhất đối nhau triệt tiêu về /bin/zsh$.

%% DANG 8D1?2-3
\textbf{Dạng 3. Tính giá trị của đa thức}
\begin{enumerate}
\item \textbf{Bước 1}: Thu gọn đa thức để biểu thức ngắn gọn, hạn chế nhầm lẫn tính toán.
\item \textbf{Bước 2}: Thay các giá trị cụ thể của các biến vào đa thức thu gọn.
\item \textbf{Bước 3}: Thực hiện phép tính lũy thừa, nhân chia rồi cộng trừ theo thứ tự.
\end{enumerate}
Mẹo: Nếu bài toán cho mối liên hệ giữa các biến (ví dụ:  + y = 2$), hãy tìm cách biến đổi đa thức xuất hiện biểu thức +y$ rồi mới thay số.
